% GENERATED FILE. Edit canonical lessons, then run npm run generate:books.
% canonical-source-hash: fnv1a64:c6b94d357d8f3084
% canonical-lessons: BN-C406-gesthaus, BN-C406-hostel, BN-C406-dhormoshala, BN-C406-thai, BN-C406-toka-deoya

\chapter{A Bed for the Night}
\label{ch:bn-a-bed-for-the-night}
\hlchaptermodality{406}

\begin{chapteropening}
\textbf{By the end of this chapter:} \emph{I can find a guesthouse, a hostel or a rest house, and somewhere to shelter.}

Complete the last lesson of chapter 406: I can find a guesthouse, a hostel or a rest house, and somewhere to shelter.
\end{chapteropening}

\section[\texorpdfstring{\bn{গেস্টহাউস} (gesṭhāus)}{গেস্টহাউস (gesṭhāus)}]{\bn{গেস্টহাউস} (gesṭhāus) — a guesthouse}

\label{lesson:BN-C406-gesthaus}



\begin{quote}
Before the new one: say the Bengali for an air conditioner, then the Bengali for to clean.
\end{quote}

\subsection*{You'll want to know: \bn{গেস্টহাউস}}
\textbf{\bn{গেস্টহাউস}} — \emph{gesṭhāus} — \textquotedblleft{}a guesthouse\textquotedblright{}.

\begin{quote}
\bn{পাহাড়ের} \bn{গেস্টহাউসটা} \bn{সস্তা} \bn{আর} \bn{পরিচ্ছন্ন।}
\end{quote}

\begin{quote}
\emph{pāhāṛer gesṭhāusṭā sôstā ār pôrichchhônnô.}
\end{quote}

\begin{quote}
The guesthouse in the hills is cheap and clean.
\end{quote}

\begin{cousinweb}
\textbf{\bn{গেস্টহাউস}} (\emph{gesṭhāus}) is borrowed from English \emph{guesthouse}.
\end{cousinweb}

\begin{grammarlens}[title={What you've built}]
One more word for getting around.
\end{grammarlens}

\subsection*{Your turn}
\begin{itemize}
\raggedright
  \item \emph{Say it:} \emph{gesṭhāus}
  \item \emph{Say it:} \emph{gesṭhāus}, once more
  \item \emph{Say it:} \emph{pāhāṛer gesṭhāusṭā sôstā ār pôrichchhônnô.} — The guesthouse in the hills is cheap and clean.
  \item \emph{Recall:} say the Bengali for an air conditioner, then the Bengali for to clean, then say \emph{gesṭhāus} again
\end{itemize}

\begin{tcolorbox}[breakable,colback=teal!4,colframe=teal!35!black,title={Before you move on}]
What does \bn{গেস্টহাউস} mean? (\textquotedblleft{}A guesthouse\textquotedblright{}.) Say it once more.
\end{tcolorbox}
\section[\texorpdfstring{\bn{হোস্টেল} (hosṭel)}{হোস্টেল (hosṭel)}]{\bn{হোস্টেল} (hosṭel) — a hostel}

\label{lesson:BN-C406-hostel}



\begin{quote}
Before the new one: say the Bengali for to clean, then the Bengali for a guesthouse.
\end{quote}

\subsection*{You'll want to know: \bn{হোস্টেল}}
\textbf{\bn{হোস্টেল}} — \emph{hosṭel} — \textquotedblleft{}a hostel\textquotedblright{}.

\begin{quote}
\bn{অরুণ} \bn{কলেজের} \bn{হোস্টেলে} \bn{থাকে।}
\end{quote}

\begin{quote}
\emph{ôrun kôlejer hosṭele thāke.}
\end{quote}

\begin{quote}
Arun lives in the college hostel.
\end{quote}

\begin{cousinweb}
\textbf{\bn{হোস্টেল}} (\emph{hosṭel}) is borrowed from English \emph{hostel}, most often a students' or workers' hall of residence.
\end{cousinweb}

\begin{grammarlens}[title={What you've built}]
One more word for getting around.
\end{grammarlens}

\subsection*{Your turn}
\begin{itemize}
\raggedright
  \item \emph{Say it:} \emph{hosṭel}
  \item \emph{Say it:} \emph{hosṭel}, once more
  \item \emph{Say it:} \emph{ôrun kôlejer hosṭele thāke.} — Arun lives in the college hostel.
  \item \emph{Recall:} say the Bengali for to clean, then the Bengali for a guesthouse, then say \emph{hosṭel} again
\end{itemize}

\begin{tcolorbox}[breakable,colback=teal!4,colframe=teal!35!black,title={Before you move on}]
What does \bn{হোস্টেল} mean? (\textquotedblleft{}A hostel\textquotedblright{}.) Say it once more.
\end{tcolorbox}
\section[\texorpdfstring{\bn{ধর্মশালা} (dhôrmôshālā)}{ধর্মশালা (dhôrmôshālā)}]{\bn{ধর্মশালা} (dhôrmôshālā) — a rest house for pilgrims}

\label{lesson:BN-C406-dhormoshala}



\begin{quote}
Before the new one: say the Bengali for a guesthouse, then the Bengali for a hostel.
\end{quote}

\subsection*{You'll want to know: \bn{ধর্মশালা}}
\textbf{\bn{ধর্মশালা}} — \emph{dhôrmôshālā} — \textquotedblleft{}a rest house for pilgrims\textquotedblright{}.

\begin{quote}
\bn{মন্দিরের} \bn{পাশের} \bn{ধর্মশালায়} \bn{আমরা} \bn{রাত} \bn{কাটালাম।}
\end{quote}

\begin{quote}
\emph{mondirer pāsher dhôrmôshālāy āmrā rāt kāṭālām.}
\end{quote}

\begin{quote}
We spent the night in the rest house beside the temple.
\end{quote}

Often free or nearly so, kept by a temple or a charitable trust.

\begin{cousinweb}
\textbf{\bn{ধর্মশালা}} (\emph{dhôrmôshālā}) is Sanskrit \emph{dharmaśālā}, \textquotedblleft{}a charitable asylum\textquotedblright{} (Monier-Williams), from \emph{dharma}, duty and merit, and \emph{śālā}, a hall: a house kept as an act of merit.
\end{cousinweb}

\begin{grammarlens}[title={What you've built}]
One more word for getting around.
\end{grammarlens}

\subsection*{Your turn}
\begin{itemize}
\raggedright
  \item \emph{Say it:} \emph{dhôrmôshālā}
  \item \emph{Say it:} \emph{dhôrmôshālā}, once more
  \item \emph{Say it:} \emph{mondirer pāsher dhôrmôshālāy āmrā rāt kāṭālām.} — We spent the night in the rest house beside the temple.
  \item \emph{Recall:} say the Bengali for a guesthouse, then the Bengali for a hostel, then say \emph{dhôrmôshālā} again
\end{itemize}

\begin{tcolorbox}[breakable,colback=teal!4,colframe=teal!35!black,title={Before you move on}]
What does \bn{ধর্মশালা} mean? (\textquotedblleft{}A rest house for pilgrims\textquotedblright{}.) Say it once more.
\end{tcolorbox}
\section[\texorpdfstring{\bn{ঠাঁই} (ṭhā̃i)}{ঠাঁই (ṭhā̃i)}]{\bn{ঠাঁই} (ṭhā̃i) — a place to stay; shelter; room}

\label{lesson:BN-C406-thai}



\begin{quote}
Before the new one: say the Bengali for a hostel, then the Bengali for a rest house for pilgrims.
\end{quote}

\subsection*{You'll want to know: \bn{ঠাঁই}}
\textbf{\bn{ঠাঁই}} — \emph{ṭhā̃i} — \textquotedblleft{}a place to stay; shelter; room\textquotedblright{}.

\begin{quote}
\bn{রাতে} \bn{গ্রামের} \bn{একটা} \bn{বাড়িতে} \bn{আমরা} \bn{ঠাঁই} \bn{পেলাম।}
\end{quote}

\begin{quote}
\emph{rāte grāmer ekṭā bāṛite āmrā ṭhā̃i pelām.}
\end{quote}

\begin{quote}
At night we found shelter in a house in the village.
\end{quote}

\bn{ঠাঁই} \bn{নেই}, there is no room: the bus was so full there was no \bn{ঠাঁই}.

\begin{cousinweb}
\textbf{\bn{ঠাঁই}} (\emph{ṭhā̃i}) is inherited: CDIAL lists Bengali \emph{ṭhā̃i}, \textquotedblleft{}a place\textquotedblright{}, under Sanskrit \emph{sth\'{\={a}}man}, \textquotedblleft{}a station, a standing place\textquotedblright{} (CDIAL 13760), from \emph{sthā}, to stand.
\end{cousinweb}

\begin{grammarlens}[title={What you've built}]
One more word for getting around.
\end{grammarlens}

\subsection*{Your turn}
\begin{itemize}
\raggedright
  \item \emph{Say it:} \emph{ṭhā̃i}
  \item \emph{Say it:} \emph{ṭhā̃i}, once more
  \item \emph{Say it:} \emph{rāte grāmer ekṭā bāṛite āmrā ṭhā̃i pelām.} — At night we found shelter in a house in the village.
  \item \emph{Recall:} say the Bengali for a hostel, then the Bengali for a rest house for pilgrims, then say \emph{ṭhā̃i} again
\end{itemize}

\begin{tcolorbox}[breakable,colback=teal!4,colframe=teal!35!black,title={Before you move on}]
What does \bn{ঠাঁই} mean? (\textquotedblleft{}A place to stay; shelter; room\textquotedblright{}.) Say it once more.
\end{tcolorbox}
\section[\texorpdfstring{\bn{টোকা} \bn{দেওয়া} (ṭokā deoyā)}{টোকা দেওয়া (ṭokā deoyā)}]{\bn{টোকা} \bn{দেওয়া} (ṭokā deoyā) — to tap; to knock lightly}

\label{lesson:BN-C406-toka-deoya}



\begin{quote}
Before the new one: say the Bengali for a rest house for pilgrims, then the Bengali for a place to stay.
\end{quote}

\subsection*{You'll want to know: \bn{টোকা} \bn{দেওয়া}}
\textbf{\bn{টোকা} \bn{দেওয়া}} — \emph{ṭokā deoyā} — \textquotedblleft{}to tap; to knock lightly\textquotedblright{}.

\begin{quote}
\bn{রাতে} \bn{কেউ} \bn{আমাদের} \bn{দরজায়} \bn{টোকা} \bn{দিল।}
\end{quote}

\begin{quote}
\emph{rāte keu āmāder dôrjāy ṭokā dilô.}
\end{quote}

\begin{quote}
In the night someone tapped on our door.
\end{quote}

Past (he/she …): \bn{তিনি} \bn{দরজায়} \bn{টোকা} \bn{দিলেন}, \emph{tini dôrjāy ṭokā dilen}.

A light tap with the knuckle; a loud knock is \bn{দরজায়} \bn{ধাক্কা}.

\begin{cousinweb}
\textbf{\bn{টোকা} \bn{দেওয়া}} (\emph{ṭokā deoyā}) joins \bn{টোকা}, \textquotedblleft{}a light tap\textquotedblright{}, to \bn{দেওয়া}, to give. \bn{টোকা} perhaps belongs with the reconstructed \emph{ṭōkk}, \textquotedblleft{}strike\textquotedblright{}, under which CDIAL lists Bengali \emph{ṭokar}, \textquotedblleft{}a rap\textquotedblright{} (CDIAL 5476); the link is hedged.
\end{cousinweb}

\begin{grammarlens}[title={What you've built}]
One more word for getting around.
\end{grammarlens}

\subsection*{Your turn}
\begin{itemize}
\raggedright
  \item \emph{Say it:} \emph{ṭokā deoyā}
  \item \emph{Say it:} \emph{ṭokā deoyā}, once more
  \item \emph{Say it:} \emph{rāte keu āmāder dôrjāy ṭokā dilô.} — In the night someone tapped on our door.
  \item \emph{Recall:} say the Bengali for a rest house for pilgrims, then the Bengali for a place to stay, then say \emph{ṭokā deoyā} again
\end{itemize}

\begin{tcolorbox}[breakable,colback=teal!4,colframe=teal!35!black,title={Before you move on}]
What does \bn{টোকা} \bn{দেওয়া} mean? (\textquotedblleft{}To tap; to knock lightly\textquotedblright{}.) Say it once more.
\end{tcolorbox}
